\expandafter\def\csname HistoricalNoteText1\endcsname{Here I always understand an imaginary quantity to mean a quantity of the form $a+b\sqrt{-1}$, as long as $b$ is not $=0$. All mathematicians of the first rank have always understood the expression in this sense. I do not think we should listen to those who have wished to call a quantity $a+b\sqrt{-1}$ imaginary only when $a=0$, but impossible when $a=0$ is not the case, since this distinction is neither necessary nor of any use. --- If imaginary quantities are to be retained in analysis at all (which, for several reasons, seems more advisable than abolishing them, provided they are established on a sufficiently sound basis), they must necessarily be regarded as just as possible as real quantities. I would therefore prefer to include real and imaginary quantities under the common name \textit{possible quantities}. In contrast, I would call a quantity \textit{impossible} if it were required to satisfy conditions that cannot be satisfied even when imaginary quantities are admitted, but in such a way that this \textit{wording} means the same as saying that no such quantity exists anywhere in the range of magnitudes. I would by no means allow a special kind of quantity to be formed from this. If someone says that a rectilinear equilateral right-angled triangle is impossible, no one will deny it. But if he wished to regard such an impossible triangle as a new kind of triangle and apply other properties of triangles to it, who could keep from laughing? This would be playing with words, or rather abusing them. --- Even the greatest mathematicians have quite often applied truths that clearly presuppose the possibility of the quantities concerned to quantities whose possibility was still doubtful. Nor would I deny that such liberties generally concern only the form and, as it were, the covering of the reasoning, which a true mathematician’s sharp eye can soon penetrate. Yet it seems more advisable, and more worthy of the eminence of a science rightly celebrated as the most perfect model of clarity and certainty, either to banish such freedoms altogether, or at least to use them more sparingly and only where even the less practised can see that the matter could have been completed without their help, perhaps less briefly, but just as rigorously. --- I would not, moreover, deny that what I have said here against the misuse of impossible quantities could in a certain respect also be objected against imaginary quantities; but I reserve their defence, together with a fuller account of this entire matter, for another occasion.}

